\subsubsection{Deduplication}

\paragraph{DISTINCT.} \texttt{DISTINCT} compares the entire select list as a
tuple. \texttt{SELECT DISTINCT a, b} returns every distinct \emph{pair}, so a
value in \texttt{a} can still repeat as long as the paired \texttt{b} differs.
It is a whole-row operator with no way to say ``dedupe on this column and keep
the rest of the row.'' It runs after \texttt{SELECT} and before
\texttt{ORDER BY}, which is why every \texttt{ORDER BY} expression must also
appear in the select list, and it treats two \texttt{NULL}s as duplicates.
\texttt{COUNT(DISTINCT col)} is a separate mechanism that dedupes the input to
one aggregate and has no effect on the rows returned.

\newpage

\paragraph{DISTINCT ON.} A PostgreSQL extension that keeps the first row of
each group sharing the same value of the parenthesized expressions, returning
all selected columns of that surviving row. ``First'' is defined by
\texttt{ORDER BY}, whose leading terms must match the \texttt{DISTINCT ON}
expressions. Omit the \texttt{ORDER BY} and the surviving row is arbitrary.

\begin{lstlisting}[style=sql, caption={Most recent close per ticker}]
SELECT DISTINCT ON (ticker)
       ticker, trade_date, close
FROM prices
ORDER BY ticker, trade_date DESC;   -- leading term must be ticker
\end{lstlisting}

\vspace{-1ex}

This is the compact form of the rank-then-filter pattern from the window
section. The \texttt{ROW\_NUMBER} equivalent needs a subquery and an outer
\texttt{WHERE rn = 1}. Which to reach for:

\begin{itemize}[label=$\circ$]
  \item \textbf{Top-1 per group in Postgres:} \texttt{DISTINCT ON} is shorter
  and generally at least as fast, since it stops at the first row of each
  group.
  \item \textbf{Top-$N$ for $N > 1$, or portability:} \texttt{ROW\_NUMBER}
  extends to \texttt{rn <= 3} and runs on any engine. \texttt{DISTINCT ON} is
  Postgres-only.
  \item \textbf{Ties:} \texttt{DISTINCT ON} always emits exactly one row per
  group, settled by the remaining \texttt{ORDER BY} terms or arbitrarily. To
  keep every tied row, use \texttt{RANK() = 1}.
\end{itemize}

Because the \texttt{ORDER BY} is dictated by the deduplication key, presenting
the result in a different order means wrapping the query in an outer
\texttt{SELECT ... ORDER BY}.
